\documentclass{ipgpmaster}
%
% Pour saisir directement les lettres accentuées :
% \usepackage[applemac]{inputenc} % pour utilisateurs mac
% \usepackage[ansinew]{inputenc}  % pour Windows ANSI
% \usepackage[latin1]{inputenc}   % pour Unix Latin1
\usepackage[T1]{fontenc}
\usepackage{amsmath}
\usepackage{float}
\usepackage{hyperref}
\hypersetup{
    colorlinks=true,
    linkcolor=blue,
    filecolor=magenta,      
    urlcolor=cyan,
    citecolor=blue,
}

\begin{document}

\checkyears

\vspace*{5mm}

%\setkeys{Gin}{draft} % images non affichées
\setkeys{Gin}{draft=false} % images affichées

%\linenumbers % Numerotation des lignes pour relecture
%Select here the language 
% \selectlanguage{french}   
\selectlanguage{english} 

\def\author{Phan Dinh Trung - 20230093}
\def\title{Robust Adaptive Face Inpainting \\ (RAFI++)}
\def\shorttitle{RAFI++: Face Mask Removal}
\def\unit{\textit{Hanoi University of Science and Technology - SOICT}}
\def\course{\textit{Project II}}
\def\spe{} % Parcours
\def\supervisor{\textit{Prof. Dr. Huynh Quyet Thang}}
\def\mydate{\today}

\Entete

\vspace*{2cm}
\begin{abstract}
% TODO: A brief summary of the problem, the proposed RAFI++ framework, key improvements over the original RAFI paper, and a summary of the results.
\end{abstract}

\newpage
\tableofcontents
\newpage

\section{Introduction}

\subsection{Background and Motivation}
Sự bùng phát của đại dịch COVID-19 và các quy định y tế sau đó đã khiến việc đeo khẩu trang trở thành một phần thiết yếu trong cuộc sống hàng ngày. Mặc dù khẩu trang đóng vai trò quan trọng trong việc bảo vệ sức khỏe cộng đồng, chúng lại che khuất một phần lớn khuôn mặt, đặc biệt là vùng mũi, miệng và cằm. Sự che khuất này gây ra những thách thức lớn đối với các hệ thống thị giác máy tính, làm giảm đáng kể hiệu suất của các ứng dụng như nhận diện khuôn mặt, phân tích cảm xúc và tương tác giữa người và máy. 

Do đó, khả năng loại bỏ khẩu trang kỹ thuật số và khôi phục lại các đặc điểm khuôn mặt bị che khuất một cách hợp lý đã trở thành một bài toán nghiên cứu vô cùng thực tiễn và có giá trị cao. Khác với các bài toán \textit{image inpainting} (phục hồi ảnh) thông thường—nơi các vùng bị khuyết có thể được lấp đầy bằng các kết cấu chung chung—việc loại bỏ khẩu trang đòi hỏi sự tuân thủ nghiêm ngặt về cấu trúc giải phẫu của khuôn mặt, tính đối xứng hai bên và khả năng duy trì bản sắc cá nhân, khiến đây trở thành một thách thức vô cùng phức tạp.

\subsection{Problem Statement}
Mục tiêu cốt lõi của bài toán này là nhận một hình ảnh khuôn mặt bị che khuất một phần bởi khẩu trang và tạo ra một khuôn mặt hoàn chỉnh, chân thực, không có khẩu trang và không chứa các chi tiết lỗi (artifacts). Đây vốn dĩ là một bài toán \textit{ill-posed} (không có nghiệm duy nhất), vì vùng bị che khuất chứa những chi tiết có tần số cao và thay đổi rất lớn giữa các cá nhân khác nhau. 

Một hệ thống phục hồi thành công không chỉ cần phân đoạn và xác định chính xác ranh giới của vùng bị che khuất, mà còn phải tổng hợp các đặc điểm giải phẫu hợp lý để có thể hòa trộn liền mạch với phần khuôn mặt phía trên không bị che khuất. Hơn nữa, các cấu trúc khuôn mặt ở phần dưới được tạo ra phải duy trì sự nhất quán về mặt ngữ nghĩa và cấu trúc với tổng thể khuôn mặt, tránh sự chuyển đổi màu sắc thiếu tự nhiên hay các đường ranh giới chắp vá dễ làm lộ ra bản chất nhân tạo của hình ảnh.

\subsection{Limitations of Previous Methods}
Các phương pháp \textit{image inpainting} ban đầu chủ yếu dựa vào các kỹ thuật truyền thống như diffusion hoặc đối sánh bản vá (patch-matching) \cite{criminisi2004region}. Dù các phương pháp này có thể xử lý tốt các kết cấu nền đơn giản, chúng lại liên tục gặp thất bại trong việc nắm bắt các thông tin ngữ nghĩa cấp cao cần thiết cho các cấu trúc khuôn mặt phức tạp. Cùng với sự phát triển của học sâu (Deep Learning), các phương pháp sử dụng Mạng đối nghịch sinh (\textit{Generative Adversarial Networks - GANs}) \cite{goodfellow2014generative} và tích chập một phần (\textit{Partial Convolutions}) \cite{liu2018image} đã cho thấy sự tiến bộ vượt bậc. Tuy nhiên, các mô hình inpainting đa dụng này thường gặp khó khăn với các ràng buộc cấu trúc cụ thể của khuôn mặt người, dẫn đến việc thường xuyên tạo ra các đặc điểm khuôn mặt bị biến dạng hoặc các vệt lỗi dễ nhận thấy xung quanh ranh giới khẩu trang.

Gần đây, một kiến trúc chuyên biệt mang tên \textit{Face Mask Removal with Region-attentive Face Inpainting} \cite{yang2024face} đã đề xuất một phương pháp gồm hai giai đoạn: một mạng phân đoạn (segmentation network) và một mạng phục hồi (restoration network). Mặc dù được thiết kế có mục tiêu rõ ràng, kiến trúc gốc này vẫn bộc lộ một số hạn chế nhất định. Mô hình chủ yếu dựa vào các cơ chế tích chập tiêu chuẩn, vốn chưa thể nắm bắt trọn vẹn các phụ thuộc ngữ cảnh tầm xa (\textit{long-range contextual dependencies}) cần thiết để tái tạo khuôn mặt nhất quán trên toàn cục. Hơn nữa, cơ chế kết hợp đặc trưng và xử lý ranh giới của nó đôi khi dẫn đến sự hòa trộn chưa tối ưu, để lại những đường nối có thể nhìn thấy bằng mắt thường giữa phần khuôn mặt phía trên và phần khuôn mặt được tổng hợp phía dưới.

\subsection{Contributions}
In this project, we propose \textit{RAFI++}, an enhanced architecture that builds upon and extends the original RAFI framework \cite{yang2024face}. First, we redesign the segmentation stage into \textit{SegNet++}, which outputs three complementary maps (soft mask, boundary map, and confidence map) instead of a single binary mask, providing the restoration stage with richer spatial guidance. Second, we improve the restoration network, now called \textit{RestoreNet++}, by introducing a decoupled channel-spatial attention module (\textit{MCSAM}), dense-style inter-block connections with learnable scaling, \textit{Skip Attention Gates} conditioned on segmentation outputs, and \textit{Residual Refine Blocks} for sharper facial details. Third, we enrich the training objective with additional \textit{SSIM}, \textit{identity}, and \textit{edge} losses, and adopt a \textit{three-stage curriculum training strategy} for more stable convergence.


\newpage
\section{Related Work}

\subsection{Deep Image Inpainting}
% TODO: General techniques (e.g., partial convolutions, gated convolutions).

\subsection{Face Mask Removal}
% TODO: Approaches specific to restoring masked faces.

\subsection{Attention Mechanisms in Vision}
% TODO: Contextual and multi-scale attention relevant to your MCSAM.


\newpage
\section{Proposed Method (RAFI++)}

\subsection{Framework Overview}
% TODO: The two-stage design (Segmentation and Restoration).

\subsection{Data Preparation Pipeline}
% TODO: Using MediaPipe Face Mesh for automatic lower-face polygon masks and boundary maps generation.

\subsection{Segmentation Network (SegNet++)}
% TODO: U-Net based Encoder-Decoder structure. Output heads: mask prediction, boundary prediction, confidence prediction.

\subsection{Restoration Network (RestoreNet++)}
% TODO: Gated Convolution Encoder, Multi-scale Channel-Spatial Attention Module (MCSAM) bottleneck, Decoder with SkipAttentionGate and ResidualRefineBlock, Soft Blending mechanism.

\subsection{Dual Discriminators}
% TODO: Image-level PatchDiscriminator, Feature-level FeaturePatchDiscriminator (VGG-based).

\subsection{Objective Functions (Losses)}
% TODO: Segmentation Losses (Mask, Boundary, Confidence), Reconstruction & Perceptual Losses (L1, SSIM, VGG Perceptual, Style, Identity, Edge), Adversarial Hinge Loss.

\subsection{Multi-Stage Curriculum Training Strategy}
% TODO: Stage 1: SegNet++ Pretraining, Stage 2: RestoreNet++ Training with frozen SegNet++, Stage 3: Joint Fine-tuning end-to-end.


\newpage
\section{Experiments}

\subsection{Dataset}
% TODO: CelebA-HQ (256x256) and train/val/test splits.

\subsection{Implementation Details}
% TODO: Hyperparameters, optimizers, hardware, and training duration.

\subsection{Evaluation Metrics}
% TODO: Segmentation: Dice score, IoU. Restoration: L1 error, PSNR, SSIM.

\subsection{Quantitative Results}
% TODO: Table comparing RAFI++ (and potentially base RAFI or other baselines if available).

\subsection{Qualitative Results}
% TODO: Visual comparison grids of masked inputs, predicted masks, and restored faces.

\subsection{Ablation Study}
% TODO: Impact of MCSAM, Impact of Dual Discriminators vs Single, Impact of the Multi-stage Curriculum.


\newpage
\section{Conclusion and Future Work}

\subsection{Conclusion}
% TODO: Summary of how RAFI++ successfully addressed the face mask removal problem.

\subsection{Future Work}
% TODO: Potential extensions (e.g., higher resolutions, handling different mask types, real-time inference).


\newpage
\begin{thebibliography}{99}
% TODO: List of cited papers (including the original RAFI paper, Gated Convolutions, VGG, MediaPipe, etc.).
\bibitem{criminisi2004region}
A. Criminisi, P. Perez, and K. Toyama, ``Region filling and object removal by exemplar-based image inpainting,'' \textit{IEEE Transactions on Image Processing}, vol. 13, no. 9, pp. 1200--1212, 2004.

\bibitem{goodfellow2014generative}
I. Goodfellow et al., ``Generative adversarial nets,'' in \textit{Advances in Neural Information Processing Systems}, 2014, pp. 2672--2680.

\bibitem{liu2018image}
G. Liu, F. A. Reda, K. J. Shih, T.-C. Wang, A. Tao, and B. Catanzaro, ``Image inpainting for irregular holes using partial convolutions,'' in \textit{Proceedings of the European Conference on Computer Vision (ECCV)}, 2018, pp. 85--100.

\bibitem{yang2024face}
M. Yang, ``Face Mask Removal with Region-attentive Face Inpainting,'' \textit{arXiv preprint arXiv:2409.06845}, 2024.
\end{thebibliography}

\newpage
\section*{Appendices}
\addcontentsline{toc}{section}{Appendices}
% TODO: Extra visual samples, detailed network configuration tables (layer by layer).

\end{document}
